\documentclass[instructions.tex]{subfiles}
\begin{document}
 
\chapter{De l’immortalité de l’ame, et de l’étendue de ses facultés.}
 
Notre âme est encore l’image de Dieu, parce qu’elle est immortelle. Les êtres matériels, les métaux, le bois, la pierre, les autres corps, sont composés de parties; ces parties peuvent s’altérer, se diviser, tomber en dissolution; mais notre âme, étant un esprit, n’a ni parties, ni mélange; elle n’a donc point de principe d’altération et de dissolution.

C’est pour cela qu’elle est incorruptible et immortelle, qu’elle est toujours la même dans sa substance, qu’elle ne prend aucun accroissement, qu’elle ne vieillit point; aussi entière et aussi parfaite dans un petit enfant que dans un homme fait, avec cette différence, qu’un enfant ayant les organes plus embarrassés, l’âme y fait des opérations plus imparfaites.

Nier l’immortalité de notre âme, c’est donner le démenti à la religion; c’est donner le démenti à Dieu, qui l’a révélée, c’est le donner à toutes les nations de la terre et à la tradition de tous les siècles, c’est enfin démentir sa propre raison: car c’est ici une de ces vérités que tout homme, qui n’a pas l’esprit et le cœur dépravés, ne peut révoquer en doute.

Disons de l’immortalité de l’ame, ce que saint Augustin a dit de la liberté, que c’est ce que la nature nous crie; et ce qui est empreint au fond de nos cœurs par le Créateur, c’est ce que tous les hommes connoissent, depuis l’école des enfans jusqu’au trône du sage Salomon; c’est ce que les bergers chantent dans les campagnes, ce que les pasteurs enseignent dans le lieu saint, ce que le genre humain annonce dans tout l’univers. On pourrait dire de ceux qui doutent de l’immortalité de l’ame, ce qu’on a dit des pyrrhoniens qui doutaient de leur existence, qu’ils sont des menteurs, qui se vantent de douter d’une chose dont il est impossible à tout homme de bon sens de douter.

Si notre âme devait mourir, l’homme ne seroit plus un ouvrage digne de Dieu. Le bien et le mal seroient également permis, puisqu’il n’y auroit ni récompense ni châtiment pour l’un ni pour l’autre; l’homme seroit sans consolation, sans frein, sans subordination; Dieu n’auroit créé les hommes que pour faire des malheureux, des scélérats, des désespérés. Il faut que le libertinage ait étrangement aveuglé un homme pour penser de la sorte.

\secondo Mais disons plus: notre âme est tellement l’image de Dieu, qu’elle est comme infinie, et participe à la puissance infinie de Dieu. Les opérations des bêtes sont bornées, elles sont matérielles et toujours les mêmes; mais l’activité de notre âme, en un sens, est sans bornes; ses opérations sont nobles et vont à l’infini.

Voyez les découvertes et les recherches de l’esprit humain dans les sciences, dans le gouvernement et dans les arts. Entrez dans ces vastes bibliothèques où l’on voit des cent volumes sur toutes les sciences imaginables. Entrez dans les académies des savans, où l’on fait des découvertes admirables sur les grandeurs de Dieu, sur ses mystères, sur les secrets et les merveilles de la nature, sur l’harmonie et le concert qui règnent dans toutes les parties de ce grand univers.

Voyez tant d’ouvrages, d’esprit, de livres d’une érudition profonde, sur la morale, sur la jurisprudence, sur l’éloquence, sur la médecine, sur l’anatomie, sur l’astronomie, sur la géométrie, et sur les différentes parties des mathématiques. Admirez, dans ces villes policées, la sagesse des lois et du gouvernement; considérez ces temples et ces édifices superbes, ornés de tout ce que la peinture et l’architecture ont de plus noble, de plus industrieux et de plus délicat.

Voyez tous ces prodigieux et immenses appareils pour la marine, pour les fortifications, pour la guerre; toutes les différentes fabriques de toute espèce, les différentes manufactures, ces ingénieuses industries de l’homme dans tous les arts, dont la variété est comme infinie.

Or, tous ces ouvrages, ces découvertes immenses dans les sciences et dans les arts, ne sont-ce pas autant d’inventions et d’opérations de l’ame et de l’esprit de l’homme? Opérations dont la moindre passe toute l’activité des animaux, opérations qui se multiplieroient et se perfectionneroient à l’infini, si la délicatesse des organes du corps pouvait seconder l’activité de l’ame. Il est donc vrai que notre âme est l’image de Dieu.

O homme! que de grandeur, que de noblesse, que de puissance sont renfermées dans vous! vous êtes, dit le Prophète, comme des dieux, vous êtes tous les enfans du Très-Haut\footnote{Dii estis, filii Excelsi omnes. Ps. 81. 6.}. Et puisque le Créateur vous a formé à son image, en vous donnant une âme si noble, capable de tant de choses, ne vous déshonorez pas, souvenez-vous de ce que vous êtes. Au lieu de vous servir de vos lumières et des facultés de votre âme, pour vous dégrader et vous perdre, que ne vous en servez-vous pour vous élever jusqu’à votre Créateur, pour l’honorer et pour vous sauver?
 
\section{Résolutions}
 
\primo J’estimerai mon âme plus que toute la terre et que tous les biens temporels du monde entier, puisqu’elle ne doit jamais mourir.

\secundo C’est pourquoi je ne négligerai plus rien pour lui procurer le ciel, bonheur qui durera autant qu’elle. 

\prayer{O mon Dieu! faites-moi la grâce de laver toutes les souillures que le péché a faites à mon âme, et de l’orner de toutes les vertus que vous voulez voir en elle pour la sauver.}
 
\end{document}
